\chapter{Context and State of the Art}
\label{ch:context}

\section{Social Infrastructure and the Accessibility Problem}
\label{sec:problem-context}

Social infrastructure, understood here as the network of preschools and schools that families use on a daily basis, differs from most other urban assets in one decisive respect: its value to a household is almost entirely determined by proximity. A kindergarten three kilometres away is, for a family without a car, functionally close to no kindergarten at all. Provision therefore cannot be evaluated by counting facilities or places at the city level. Two cities with identical numbers of places per thousand residents can offer radically different lived experience depending on how those places are distributed relative to where people actually live.

This distinction between aggregate provision and spatial accessibility is the starting point of the present work. A city that grows rapidly at its periphery, as capitals under sustained public investment typically do, accumulates a characteristic pattern: an inherited core with dense and often underused facility stock, and newly built residential districts where population arrives years before the corresponding social infrastructure. The mismatch is not primarily a shortage of investment. It is a coordination problem between the timing and location of housing construction and the timing and location of facility construction.

Astana illustrates the pattern. The capital has benefited from substantial public investment and broad access to social services, while its construction boom has placed intensive demands on land and financial resources \cite{astana2025}. The resulting pressure on schools is officially recognised: the national project Comfortable School records that the student-place deficit is most acute in five regions, among them Astana, which together account for 59 per cent of the national deficit, and attributes the persistence of three-shift schooling to the absence of quality planning of school construction that accounts for the place deficit and the growth of the child population \cite{comfortschool2022}.

Quantifying that mismatch requires three things that this chapter reviews in turn: a defensible measure of accessibility that accounts for both supply capacity and demand competition, a way of summarising how unequally that measure is distributed across the population, and a method for identifying where new capacity would do the most good. A fourth question, less often asked, concerns the planner rather than the plan: given that siting decisions are already being made, what considerations do they actually reflect, and how far do they diverge from what an accessibility-driven allocation would suggest.

\section{Measuring Spatial Accessibility}
\label{sec:accessibility-measures}

\subsection{Container and Buffer Measures}

Accessibility has been understood since Hansen as the potential of opportunities for interaction \cite{hansen1959}, and the measures that operationalise it are commonly grouped into infrastructure-based, location-based, person-based and utility-based perspectives \cite{geurs2004}. The measures relevant to facility planning are location-based, and they range from simple to demanding.

The simplest accessibility measures partition space administratively and compute a supply-to-demand ratio within each partition, for example places per thousand children in a district. These container measures are attractive because they use only data that administrations already collect, but they carry a defect that is fatal for fine-grained analysis: they assume that residents use facilities inside their own administrative unit and no others. A household living fifty metres from a district boundary is treated as having no access to the school on the other side of the street. The measure is also sensitive to the arbitrary geometry of the partition, an instance of the modifiable areal unit problem.

Buffer or isochrone measures improve on this by defining accessibility relative to the household rather than the district: a location is served if a facility lies within a fixed distance or travel time. This is the logic embedded in most national planning standards, which express provision requirements as service radii. Buffer measures are transparent and map directly onto regulatory compliance, which makes them indispensable for the normative part of any assessment. Their limitation is that they are binary and capacity-blind. A location within 300 metres of a single overcrowded kindergarten and a location within 300 metres of three half-empty ones receive the same score. The choice between such measures is not neutral: whether a public service is judged equitably distributed can depend on the accessibility measure used to assess it \cite{talen1998}.

\subsection{The Floating Catchment Area Family}

The floating catchment area approach was developed to combine the household-centred logic of buffers with the supply-and-demand logic of container measures. Radke and Mu \cite{radke2000} introduced the two-stage formulation, and Luo and Wang \cite{luo2003} popularised it as the two-step floating catchment area (2SFCA) method. The method proceeds in two passes. In the first, a catchment is drawn around each supply location, the demand falling inside it is summed, and a supply-to-demand ratio is computed for that facility. In the second, a catchment is drawn around each demand location and the ratios of all facilities falling inside it are summed. The result is an accessibility score per demand location that reflects both how much capacity is reachable and how many other people compete for it.

The original formulation applies a uniform weight inside the catchment and zero outside, which reintroduces a discontinuity at the boundary and treats a facility at the edge of the catchment as equivalent to one next door. Luo and Qi \cite{luo2009} addressed this with the enhanced two-step floating catchment area (E2SFCA) method, which divides the catchment into travel-time zones and applies a decreasing weight to each. Subsequent variants have refined the treatment further. Wang \cite{wang2012} reviewed the family and formalised the connection between the floating catchment approach and gravity-based accessibility models, showing that 2SFCA is a special case of a gravity model with a step impedance function. Delamater \cite{delamater2013} observed that the standard formulation implicitly assumes an optimally configured supply system and proposed a modified metric (M2SFCA) that penalises configurations in which capacity is poorly placed relative to demand. Wan, Zou and Sternberg \cite{wan2012} added a competition step that apportions demand between nearby facilities, in order to reduce the overestimation of demand inherent in the two-step formulation.

Recent applications have extended the family in two directions relevant to the present work. The first is multi-modal formulation, in which the impedance is computed separately for walking, public transport and driving and combined according to observed or assumed mode shares. The second is the coupling of the accessibility measure with inequality statistics and spatial autocorrelation diagnostics within a single analytical workflow, an approach exemplified by recent studies of elderly care provision \cite{eldercare2025}. This combination, accessibility measurement followed by equity diagnosis followed by cluster detection, has become close to a standard template. The method has also been carried over from healthcare to education: it has been used to measure access to primary schools in Chengdu, with a multi-modal formulation and a modelled distribution of primary school-age children \cite{chengdu2024}, and access to upper secondary education in an Italian region, where geographical distance proved to be a barrier for rural students that neither economic status nor digital infrastructure compensated for \cite{italy2022}.

\subsection{Impedance Functions and Catchment Thresholds}

Two parameter choices dominate the behaviour of any floating catchment model: the shape of the distance-decay function and the size of the catchment. Neither is determined by the data. Common choices for the decay function include step, linear, Gaussian, exponential and power forms, and the literature offers no consensus, in part because the appropriate shape depends on the service. Attendance at a preschool falls off far more sharply with distance than attendance at a specialised clinic. McGrail \cite{mcgrail2012} evaluated the two main groups of proposed refinements, distance decay within the catchment and variable catchment size, and compared how they behave in metropolitan and rural populations.

Impedance is not constant through the year either. On the best practically obtainable winter pavement an average pedestrian needs about one minute per kilometre longer than on bare pavement \cite{winterwalk2021}, and weather conditions affect walking speed in severely cold climates \cite{coldwalk2020}. For a city with a long continental winter this is a reason to compute a seasonal variant of the travel-time matrix rather than a single annual one.

The catchment threshold has an analogous problem, with one important difference in the context of this work. For services governed by a national planning standard, the threshold is not a free parameter to be tuned but a regulatory quantity to be respected. Anchoring the catchment on the statutory service radius rather than on a value borrowed from a foreign study makes the resulting accessibility surface directly interpretable as a compliance measure, and removes one degree of analyst discretion. Where the standard is silent, the sensitivity of the results to the remaining choices must be reported rather than assumed away; this obligation is taken up in the validation design of Chapter~\ref{ch:results}.

\section{Equity Metrics for Accessibility Distributions}
\label{sec:equity}

An accessibility surface is a map. To make claims about fairness it must be reduced to statistics that describe how unequally accessibility is distributed across people rather than across space. The distinction matters: an accessibility deficit in an empty industrial zone is not an equity problem, whereas the same deficit in a dense residential district is.

The standard instruments are borrowed from income inequality analysis. The Lorenz curve plots the cumulative share of accessibility against the cumulative share of population, ordered from least to most advantaged, and the Gini coefficient summarises the area between that curve and the line of perfect equality. The same instruments are applied to public services, including public transport \cite{delbosc2011} and healthcare \cite{neutens2015}, and the interpretation carries over directly: a population-weighted Gini of zero would mean every resident enjoys identical accessibility.

The Gini coefficient is a scalar and therefore says nothing about where inequality is located. The Theil index \cite{theil1967} is preferred where that question matters, because it decomposes additively into a within-group and a between-group component. Applied with administrative districts as the grouping variable, the decomposition separates inequality that exists between districts, which is a matter for city-level investment allocation, from inequality inside districts, which is a matter for local siting.

A recognised limitation of the whole family is that these indices are aggregate and group-blind. A distribution can appear acceptably equal overall while a specific and identifiable subpopulation is systematically disadvantaged. Evaluating the equity of accessibility therefore requires attention to how accessibility is distributed over particular areas and population groups, not only to its aggregate spread \cite{lucas2016}. In the present work this limitation is addressed not by replacing the indices but by supplementing them with two additional views: the share of each demographic cohort falling outside the statutory service radius, which is a direct and legible measure of vertical equity against a legal benchmark, and spatial cluster detection, which locates disadvantage rather than merely quantifying it.

Spatial cluster detection relies on measures of spatial autocorrelation. The global Moran's~$I$ statistic \cite{moran1950} tests whether the accessibility surface is more spatially clustered than would be expected under spatial randomness. Because a single global statistic can conceal opposing local patterns, it is complemented by local indicators of spatial association \cite{anselin1995}, which classify each unit according to whether it and its neighbours are jointly high, jointly low, or discordant. The jointly-low clusters are the deficit areas that a siting intervention should target, and identifying them statistically rather than by visual inspection of a choropleth map is what distinguishes a defensible diagnosis from an impressionistic one.

\section{Facility Location Optimization}
\label{sec:location-models}

Where accessibility measurement is descriptive, facility location modelling is prescriptive. The field originates with Hakimi \cite{hakimi1964}, who established that for the $p$-median problem on a network, an optimal solution exists at the nodes of the graph, reducing a continuous problem to a finite combinatorial one. ReVelle and Swain \cite{revelle1970} gave the integer programming formulation of this central facilities problem that remains its standard statement. The $p$-median model selects $p$ facility locations so as to minimise the total demand-weighted travel distance, and is the natural formulation when the planning objective is average convenience.

Church and ReVelle \cite{church1974} introduced the maximal covering location problem (MCLP), which changes the objective: rather than minimising average distance, it maximises the amount of demand falling within a fixed coverage radius, subject to a budget of $p$ facilities. The distinction is consequential. The $p$-median model will happily improve the situation of already well-served residents if doing so lowers the weighted average, whereas MCLP is indifferent to improvements inside the covered set and rewards only the conversion of uncovered demand into covered demand. Where a planning standard defines a service radius, MCLP is the formulation that matches the regulatory objective. Church and ReVelle \cite{church1976} later established the theoretical and computational relationships between the two models and the location set covering problem, and Murray \cite{murray2016} reviewed the four decades of extensions that followed, including capacitated, multi-period, hierarchical and probabilistic variants. Broader reviews of strategic facility location \cite{owen1998} and of covering problems in particular \cite{farahani2012} document how widely the models have been applied. Both models, together with a capacitated coverage variant, are now available in mainstream geographic information system software \cite{esrilocalloc}, which has made them a routine planning instrument rather than a research curiosity.

Two extensions matter for social infrastructure specifically. The first is capacity. An uncapacitated model may assign more demand to a facility than it can serve; since school and preschool places are the binding constraint in practice, a capacitated formulation is required whenever capacity data are available, and capacitated versions of both the set covering and the maximal covering problems exist for exactly this reason \cite{current1988}. Capacitated location modelling has been applied to public primary schools directly, first to evaluate an existing network by identifying areas of shortage and excess, then to propose relocation \cite{pizzolato2004}. The second is the treatment of competing objectives. Coverage and equity do not move together: a solution that maximises total covered demand will tend to place facilities in dense areas, which may leave a sparse but severely underserved periphery untouched. Collapsing the two objectives into a single weighted sum requires choosing weights, and the choice is difficult to justify. Reporting a Pareto frontier instead, obtained either by the epsilon-constraint method \cite{mavrotas2009} or by an evolutionary multi-objective algorithm such as NSGA-II \cite{deb2002}, presents the trade-off to the decision maker rather than resolving it on their behalf, and removes the most common line of criticism directed at multi-criteria siting studies.

\section{Machine Learning and Explainability in Urban Analytics}
\label{sec:ml-urban}

Gradient-boosted decision tree ensembles are widely used for tabular data, in implementations such as XGBoost \cite{chen2016}, LightGBM \cite{ke2017} and CatBoost \cite{prokhorenkova2018}. Their appeal in this domain is practical: they handle heterogeneous features and non-linear interactions without manual specification, tolerate missing values, and train quickly enough to permit thorough cross-validation.

Because the resulting models are not directly interpretable, their adoption in planning contexts has been accompanied by post-hoc explanation methods, most prominently SHAP \cite{lundberg2017}, which attributes a prediction to its input features on the basis of Shapley values from cooperative game theory, with an efficient exact algorithm available for tree ensembles \cite{lundberg2020}. For geographical data specifically, SHAP-explained gradient boosting has been shown in simulation to recover spatial effects similar to those estimated by spatial lag and multiscale geographically weighted regression models \cite{li2022}, and Shapley values have since been extended to measure the effect of location itself within a machine learning model \cite{li2024}. SHAP therefore serves two purposes here: ranking drivers globally and interrogating individual locations.

Applications of machine learning to facility siting specifically are also established. A recent study in Tianjin built a suitability model for primary schools on a 500\,m grid from the existing distribution of schools, using random forest and XGBoost; it reported evaluation metrics above 85 per cent that exceeded logistic regression, support vector machine, nearest-neighbour and single decision tree baselines, and used SHAP to identify population density, proximity to other educational institutions and transport accessibility as the main drivers of suitability \cite{tianjin2025}. Related studies combine machine learning with multi-criteria decision analysis and feature ranking into hybrid siting frameworks \cite{schoolsiting2025}. It is therefore not accurate to present machine learning for school siting as novel in itself, and this work does not do so.

What remains open is the interpretation of what such a model represents. A suitability model trained on the locations of existing facilities does not learn where facilities \emph{should} be. It learns where they \emph{were placed}, and therefore encodes the historical preferences, constraints and possibly the errors of the planners who placed them. Read as a normative instrument, this is a defect and a mechanism for perpetuating existing inequity. Read as a descriptive instrument, it is informative in a different way: it becomes a revealed-preference model of the planning process, and its SHAP attribution becomes evidence about what siting decisions actually optimise. Setting that descriptive model against a normative optimisation over the same demand surface converts the ambiguity into the central research question rather than leaving it as an unexamined assumption. This framing, rather than the use of machine learning as such, is where the present work departs from the studies cited above.

\section{Validation Under Spatial Dependence}
\label{sec:spatial-validation}

Any predictive model fitted to spatial data faces a validation problem that is easy to overlook and severe in its consequences. Nearby locations resemble one another, so a randomly selected test set contains observations that are near-duplicates of observations in the training set. The model is then evaluated on data it has effectively already seen, and the reported performance is optimistically biased. Roberts et al. \cite{roberts2017} set out the general case for structured cross-validation in the presence of temporal, spatial, hierarchical and phylogenetic dependence, and the finding has since been replicated across geospatial modelling: spatially autocorrelated training and validation samples inflate performance estimates and misrepresent how a model will generalise across its intended application domain. In a large-scale forest biomass mapping exercise, for example, a random forest that validated well under a standard non-spatial procedure showed poor predictive power once spatial autocorrelation was taken into account \cite{ploton2020}.

The remedy is to construct folds that are spatially separated, either by partitioning the study area into contiguous blocks or by buffering each test observation against training observations within the range of autocorrelation. Neither is free of difficulty. Blocking reduces the effective training set and can itself bias estimates in the opposite direction, and there is an active methodological debate about which strategy is appropriate for which prediction task, including the argument that spatial cross-validation is not the right way to evaluate map accuracy \cite{wadoux2021}. What is not in dispute is that random $k$-fold cross-validation on spatially dependent data reports a number that does not mean what it appears to mean.

This has a direct bearing on the siting literature reviewed in Section~\ref{sec:ml-urban}. Reported accuracies above eighty-five per cent for facility suitability models are difficult to interpret without knowing how the folds were constructed, since the target variable in such models, the presence of a facility, is strongly spatially clustered by construction. The present work therefore treats the validation protocol as a substantive component of the methodology rather than as a routine step, adopts spatially blocked folds, and reports Moran's~$I$ on model residuals as evidence that the blocking achieved what it was intended to achieve.

\section{Existing Decision Support Platforms}
\label{sec:existing-systems}

Software that supports planning decisions with spatial analysis is usually grouped under the heading of planning support systems, a varied family that ranges from general scenario sketching tools to dedicated simulation models \cite{geertman2004}. A comparative assessment of the most relevant representatives shows what an open research framework has to add and what it should not attempt.

UrbanSim simulates urban development for land use, transportation and environmental planning \cite{waddell2002}. Its strength is a behavioural model of how households, businesses and real estate development respond to policy and infrastructure over time at the metropolitan scale. Its cost is a substantial data and calibration effort, and its outputs are forecasts of development patterns rather than siting recommendations for a specific facility class.

CommunityViz is an extension to ArcGIS for scenario, suitability, build-out and impact analysis, offering a library of built-in indicator models that planners can combine or replace with their own \cite{communityviz}. It is a general-purpose indicator and scenario environment rather than a dedicated workflow for accessibility measurement, equity diagnosis and facility location.

Esri ArcGIS provides a location-allocation solver whose problem types include the $p$-median problem, maximal coverage and a capacitated coverage variant \cite{esrilocalloc}. These are generic solvers: they take demand points, candidate facilities and an impedance as given, and leave the estimation of demand, the measurement of accessibility, the diagnosis of inequity and any connection to a national planning standard to the analyst.

Three gaps follow from this comparison and define the scope of the present framework. First, none of these tools is anchored to the Kazakh planning standard, so compliance against statutory service radii is not a native output. Second, none couples a descriptive model of observed siting behaviour with a normative optimisation, which is the central analytical contribution proposed here. Third, CommunityViz and ArcGIS are licensed products and UrbanSim requires a modelling effort far beyond the question addressed here, whereas the reproducibility requirements of academic work favour an open, containerised implementation. The framework developed in this work is accordingly narrower in scope than any of the platforms reviewed, and deeper in the specific analytical chain from population grid to accessibility to equity diagnosis to optimised siting recommendation.

\section{Regulatory and Institutional Framework}
\label{sec:regulatory}

Three instruments set the parameters that a study of school and preschool provision in Kazakhstan has to respect.

The first is the national planning standard SP RK 3.01-101-2013 \cite{sprk2013}. Its table of service radii for facilities placed in residential areas sets 300\,m for preschools in cities, 500\,m for preschools in rural settlements and small towns with low-rise housing, and 500\,m for general schools. The radii may be increased by up to 10 per cent in settlements with a complex layout or under reconstruction, and schools may instead be reached by transport within 15 minutes for primary grades and 50 minutes for the rest. The same standard sets provision norms: preschool places for up to 85 per cent of children, or up to 180 places per thousand residents in new-build areas without demographic data, and school places for all pupils of grades 1 to 9 and up to 75 per cent of grades 10 and 11 in a single shift, or at least 180 places per thousand residents in new-build areas. It also prescribes plot sizes, from 43\,m$^2$ per pupil for small schools down to 19\,m$^2$ per pupil for schools of 900 to 1,200 pupils, which puts the plot of a 2,000-seat school at about 3.8\,ha.

The second is the system of regional standards for settlements approved by government decree in 2025 \cite{decree7102025}. It states the principle of twenty-minute walking access to services and, for the capital and the cities of republican significance, refers the calculation of school and preschool places back to the planning standard.

The third is the national project Comfortable School \cite{comfortschool2022}, the instrument through which most new school capacity in Astana has been delivered since 2023, and whose sites are used in this work as a test of siting decisions.

The standard does not state whether a service radius is measured in a straight line or along the walking network. Planning practice usually applies it as a straight-line distance; a child, however, walks along streets. Both readings are reported in Chapter~\ref{ch:results}, and the difference between them proves to be large.

\section{Research Gap and Problem Statement}
\label{sec:gap}

The methodological components reviewed above are individually mature. Floating catchment accessibility measurement, inequality decomposition, spatial cluster detection, covering and median location models, and explainable gradient boosting are each supported by an established literature and by working software. Their combination into a single workflow is also no longer novel; the accessibility-then-equity-then-clusters template appears repeatedly in recent applications.

Three gaps nevertheless remain, and together they define this work.

The first is geographic. Scholarship on Astana has examined the politics and lived experience of its state-led modernisation \cite{koch2014, laszczkowski2016} and, more recently, the resource and environmental costs of its infrastructure programmes \cite{astana2025}. A fine-grained quantitative assessment of access to its social infrastructure was not found in the literature reviewed for this work. Kazakhstan's capital is an informative case precisely because it is atypical: sustained public investment, rapid peripheral residential growth and an active national school construction programme produce a fast-moving supply and demand configuration that most accessibility studies, conducted in slower-changing cities, do not capture.

The second is methodological, and concerns the interpretation of learned siting models. The literature contains models that predict facility suitability from the observed distribution of facilities, but the resulting attributions are generally presented as guidance for future siting rather than as evidence about past decision making. The two readings have opposite implications, and distinguishing between them requires an independent normative benchmark over the same demand surface. Constructing that benchmark, and measuring the divergence between what was built and what a coverage- or equity-optimal allocation would have selected, is the analytical core of this work.

The third is evaluative. Siting studies are difficult to validate because the counterfactual is unobservable: there is normally no way to know whether a recommended location would have been better than the one chosen. The Astana case admits an unusually direct test. A national school construction programme \cite{comfortschool2022} has delivered a known set of facilities at known locations with known capacities over a known period, which makes it possible to reconstruct the pre-programme state, run the optimisation with the realised number of facilities as the budget, and compare the recommended set against the realised set. Combined with explicit baselines and a spatially blocked validation protocol, this converts an exercise that is usually assessed on plausibility into one assessed on measured agreement.

\subsection{Aim and Objectives}

The aim of this work is to develop and validate a data-driven framework that assesses spatial accessibility to social infrastructure, diagnoses its inequity and derives siting recommendations, and to apply it to schools and preschool organisations in Astana using only openly available data.

The objectives are:

\begin{enumerate}[nosep]
  \item to assemble from open sources an inventory of schools and preschool organisations and a population surface below the district level, and to measure the accuracy of both against independent references;
  \item to compute an enhanced two-step floating catchment area index of accessibility for each facility type, parameterised by the national planning standard;
  \item to describe how that accessibility is distributed, using population-weighted inequality indices, their decomposition by district, statutory compliance rates and statistically identified deficit clusters;
  \item to model the observed siting decisions of the national school construction programme under a spatially blocked validation protocol and interpret the fitted model;
  \item to solve covering and median location models for the same number of schools and compare the recommended allocation with the realised one;
  \item to test the stability of the conclusions under alternative population surfaces, supply measures and distance decay functions.
\end{enumerate}

\subsection{Object and Subject of Research}

The object of research is the spatial configuration of school and preschool provision and of residential demand for it in Astana. The subject of research is the set of data-driven methods by which accessibility to that provision is measured, its inequity diagnosed, its observed siting explained and its future siting optimised.

\section{Research Hypotheses}
\label{sec:hypotheses}

The work tests four hypotheses. Each is stated so that the data can contradict it.

\begin{enumerate}[label=H\arabic*.,nosep,leftmargin=*]
  \item Population below the district level can be estimated from open sources more accurately than by the standard gridded product: the number of polling-station territories reproduces the official population of districts more closely than WorldPop and than residential floor area.
  \item Accessibility to state schools and preschool organisations is unevenly distributed across Astana, and the largest deficits per child lie in the districts of recent residential growth rather than in the established centre.
  \item The spatial pattern of deficits found under H2 persists when the population surface, the supply measure and the distance decay function are varied within their plausible ranges.
  \item The sites chosen by the national school construction programme are better predicted by pre-programme accessibility deficits and urban form than by population density alone, and an allocation that maximises coverage of underserved children with the same number of schools would have covered more of them than the realised programme.
\end{enumerate}

H1 concerns the measurement on which every later result depends and is tested first, in Chapter~\ref{ch:methods}. H2 and H3 are tested in Chapter~\ref{ch:results} on the accessibility surface. H4 is tested by backtesting against the realised programme.
